\begin{afterword}
\summaryhead

The post continues the story of the young author's mistake. In 1996 Eliezer1996, primed by admiration of intelligence, by technophilia and by the broad sympathies of science fiction, set out to build superintelligence and assumed it would know what is right. On a transhumanist mailing list people objected: morality is arbitrary, a superintelligence would seize everything, humans should come first, an AI needs a control system. Eliezer1996 found all four obviously wrong and refuted them.

The author explains why winning felt like being right. Traditional Rationality pictures science as a fair fight, and finding flaws in an argument is easier than finding the truth. A passage from the author's ``Twelve Virtues of Rationality'' warns that beginners win arguments and overrate themselves. The lesson is that no one else will argue you out of your mistakes, and that one should take Nature, not people, as the opponent. That only a full reduction of morality could have persuaded Eliezer1996 measures the young author's level; the philosophers on the list did not help. The next stage, perhaps in 1997, is left for the next post.

\responsehead

As memoir, the post matches the record. The Extropians mailing list archive for November and December 1996 has the young author's confidence (``The Powers will be ethical''), objections of the kinds the post recalls, and quick, emphatic replies. It also has the ``personal philosophy'' of looking for flaws, stated on 2 December 1996.

The lesson is not new to the book, and the post acknowledges it: it quotes the author's own ``Twelve Virtues of Rationality'' for it. What the post adds is the case.

The archive, read beside the author's later writing, adds something the post leaves implicit. Two of the objections the young author beat reached conclusions close to the adult author's view in 2008: that no argument moves every possible mind, and that an unfriendly AI would treat us as raw material. The young author won those arguments, and their conclusions were in part right. That fits the post's lesson that winning an argument is not the same as being right.

In short: a memoir that the 1996 mailing-list archive supports, illustrating a lesson from the author's own ``Twelve Virtues of Rationality.''
\end{afterword}
